\chapter{2005年山东卷（理）}
\section{单选题}
\begin{problem}
\(\frac{1 - \frac{1}{1 - x}}{(1 + 1)^2}\) 1+1=（\quad）

\choices
    {i}
    {-i}
    {1}
    {\(-1\)}
\end{problem}

\begin{problem}
函数 \( y = \frac{1 - x}{x} (x \neq 0) \) 的反函数图象大致是（\quad）

\choices
    {}
    {}
    {}
    {}
\end{problem}

\begin{problem}
已知函数 \( y = \sin \left( x - \frac{\pi}{12} \right) \cos \left( x - \frac{\pi}{12} \right) \)，则下列判断正确的是（\quad）

\choices
    {此函数的最小周期为 \(2\pi\)，其图象的一个对称中心是 \(\left(\frac{\pi}{12}, 0\right)\)}
    {此函数的最小周期为 \(\pi\)，其图象的一个对称中心是 \(\left(\frac{\pi}{12}, 0\right)\)}
    {此函数的最小周期为 \(2\pi\)，其图象的一个对称中心是 \(\left(\frac{\pi}{6}, 0\right)\)}
    {此函数的最小周期为 \(\pi\)，其图象的一个对称中心是 \(\left(\frac{\pi}{6}, 0\right)\)}
\end{problem}

\begin{problem}
下列函数既是奇函数，又在区间 \([-1, 1]\) 上单调递减的是（\quad）

\choices
    {\(f(x) = \sin x\)}
    {\(f(x) = -|x + 1|\)}
    {\(f(x) = \frac{1}{2} (a^x + a^{-x})\)}
    {\(f(x) = \ln \frac{2 - x}{2 + x}\)}
\end{problem}

\begin{problem}
如果 \(\left(3x - \frac{1}{\sqrt[3]{x^2}}\right)^n\) 的展开式中各项系数之和为 128，则展开式中 \(\frac{1}{x^3}\) 的系数是（\quad）

\choices
    {7}
    {\(-7\)}
    {21}
    {\(-21\)}
\end{problem}

\begin{problem}
函数 \( f(x) = \begin{cases} \sin(\pi x^2), & -1 < x < 0, \\ \e^{x - 1}, & x \geqslant 0. \end{cases} \) 若 \( f(1) + f(\e) = 2 \)，则 \( \alpha \) 的所有可能值为（\quad）

\choices
    {1}
    {\(-\frac{\sqrt{2}}{2}\)}
    {\(1, -\frac{\sqrt{2}}{2}\)}
    {\(1, \frac{\sqrt{2}}{2}\)}
\end{problem}

\begin{problem}
已知向量 \(\bs{a}, \bs{b}\)，且 \(\overrightarrow{AB} = \bs{a} + 2\bs{b}, \overrightarrow{BC} = -5\bs{a} + 6\bs{b}, \overrightarrow{CD} = 7\bs{a} - 2\bs{b}\)，则三、解答题 一定共线的三点是（\quad）

\choices
    {\(A, B, D\)}
    {\(A, B, C\)}
    {\(B, C, D\)}
    {\(A, C, D\)}
\end{problem}

\begin{problem}
设地球的半径为 \(R\)，若甲地位于北纬 \(45^{\circ}\) 东经 \(120^{\circ}\)，乙地位于南纬 \(75^{\circ}\) 东经 \(120^{\circ}\)，则甲、乙两地的球面距离为（\quad）

\choices
    {\(\sqrt{3} R\)}
    {\(\frac{\pi}{6} R\)}
    {\(\frac{5\pi}{6} R\)}
    {\(\frac{2\pi}{3} R\)}
\end{problem}

\begin{problem}
10 张娄穷中只有 3 张有效，5 个人购买，每人 1 张，至少有 1 人中奖的概率是（\quad）

\choices
    {\(\frac{3}{10}\)}
    {\(\frac{1}{12}\)}
    {\( \frac{1}{2} \)}
    {\(\frac{11}{12}\)}
\end{problem}

\begin{problem}
设集合 \(A\)、\(B\) 是全集 \(U\) 的两个子集，则 \(A \subsetneq B\) 是 \((\complement_U A) \cup B = U\) 的（\quad）

\choices
    {充分不必要条件}
    {必要不充分条件}
    {充要条件}
    {既不充分也不必要条件}
\end{problem}

\begin{problem}
\(0 < \alpha < 1\) ，下列不等式一定成立的是（\quad）

\choices
    {\(|\log_{1 + \alpha}|(1 - \alpha)| + |\log_{1 + \alpha}|(1 + \alpha)| > 2\)}
    {\(|\log_{1 + \alpha}|(1 - \alpha)| < |\log_{1 - \alpha}|(1 + \alpha)|\)}
    {\(|\log_{1 + \alpha}|(1 - \alpha) + \log_{1 + \alpha}|(1 + \alpha)| < |\log_{1 + \alpha}|(1 - \alpha)| + |\log_{1 - \alpha}|(1 + \alpha)|\)}
    {\(|\log_{1 + \alpha}|(1 - \alpha) - \log_{1 + \alpha}|(1 + \alpha)| < |\log_{1 + \alpha}|(1 - \alpha)| - |\log_{1 - \alpha}|(1 + \alpha)|\)}
\end{problem}

\begin{problem}
设直线 \(k: 2x + y + 2 = 0\) 关于原点对称的直线为 \(l'\)，若 \(l'\) 与椭圆 \(x^2 + \frac{y^2}{4} = 1\) 的交点为 \(A, B\)，点 \(P\) 为椭圆上的动点，则使 \(\triangle PAB\) 的面积为 \(\frac{1}{2}\) 的点 \(P\) 的个数为（\quad）

\choices
    {1}
    {2}
    {3}
    {4}
\end{problem}

\section{填空题}
\begin{problem}
\(\lim_{n\to \infty}\frac{C_n^2 + 2C_{n-1}^n - 2}{(n + 1)^2} =\)
\end{problem}

\begin{problem}
设双曲线 \( \frac{x^{2}}{a^{2}} - \frac{y^{2}}{b^{2}} = 1 (a > 0, b > 0) \) 的右焦点为 F，右准线 l 与两条渐近线交于 P，Q 两点，如果 \( \triangle PQF \) 是直角三角形，则双曲线的离心率 e =\fillinblank{} \fillinblank{}.
\end{problem}

\begin{problem}
设 \(x\)、\(y\) 满足约束条件 \(\left\{ \begin{array}{l} x + y \leqslant 5, \\ 3x + 2y \leqslant 12 \\ 0 \leqslant x \leqslant 3, \\ \text{大的点} (x, y) \text{是} \end{array} \right.\) 则使得目标函数 \(z = 6x + 5y\) 的最
\end{problem}

\begin{problem}
已知 \(m\)、\(n\) 是不可用的直线，\(\alpha\)、\(\beta\) 是不重合的平面，给出下列命题: \circled{1} 若 \(\alpha // \beta, m \subset \alpha, n \subset \beta\)，则 \(m // n\); \circled{2} 若 \(m, n \subset \alpha, m \nmid \beta, n \nmid \beta\)，则 \(\alpha \nmid \beta\); \circled{3} 若 \(m \perp \alpha, n \perp \beta, m \nmid n\)，则 \(\alpha \nmid \beta\); \circled{4} \(m, n\) 是两条异面直线，若 \(m \cap \alpha, m \nmid \beta, n \cap \alpha, n \nmid \beta\)，则 \(\alpha \nmid \beta\). 上面的命题中，真命题的序号是(写出所有真命题的序号)
\end{problem}

\begin{problem}
已知向量 \(\bs{m} = (\cos \theta, \sin \theta)\) 和 \(\bs{n} = (\sqrt{2} - \sin \theta, \cos \theta), \theta \in (\pi, 2\pi)\)，且 \(|\bs{m} + \bs{n}| = \frac{8\sqrt{2}}{9}\)，求 \(\cos \left(\frac{\theta}{2} + \frac{\pi}{8}\right)\) 的值 \fillinblank{}.
\end{problem}

\begin{problem}
袋中装有黑球和白球共7个，从中任取2个球都是白球的概率为 \(\frac{1}{3}\)，用甲、乙两人从袋中轮流摸取1球，甲先取.乙后取，然后甲再取…，取后不放回，直到两人中有一人取到白球时既终止.每个球在每一次被取出的机会是等可能的.用《表示取球终止所需要的取球次数. (1) 求袋中所有白球的个数; (2) 求随机变量 \( \zeta \) 的概率分布; (3) 求甲取到白球的概率 \fillinblank{}.
\end{problem}

\begin{problem}
已知 \(x = 1\) 是函数 \(f(x) = mx^3 - 3(m + 1)x^2 + nx + 1\) 的一个极值点，其中 \(m, n \in \R, m < 0\). (1) 求 \(m\) 与 \(n\) 的关系式; (2) 求 \( f(x) \) 的单调区间; (3) 当 \(x \in [-1, 1]\) 时，函数 \(y = f(x)\) 的图象上任意一点的切线斜率恒大于 \(3m\)，求 \(m\) 的取值范围 \fillinblank{}.
\end{problem}

\begin{problem}
已知数列 \(\{a_{n}\}\) 的首项 \(a_1 = 5\)，而 \(n\) 项和为 \(S_{n}\)，且 \(S_{n+1} = 2S_{n} + n + 5 (n \in \mathbf{N}^{*})\). (1) 证明数列 \(\{a_{n} + 1\}\) 是等比数列; (2) 令 \( f(x) = a_1 x + a_2 x^2 + \cdots + a_n x^n \)，求函数 \( f(x) \) 在点 \( x = 1 \) 处的导数 \( f'(1) \) 并比较 \( 2f'(1) \) 与 \( 23n^2 - 13n \) 的大小 \fillinblank{}.
\end{problem}

\begin{problem}
已知动圆过定点 \(\left(\frac{p}{q}, 0\right)\)，且与直线 \(x = -\frac{p}{2}\) 相切，其中 \(p > 0\). (1) 求动圆圆 \(C\) 的轨迹的方程. (2) 设 \(A\)、\(B\) 是轨迹 \(C\) 上异于原点 \(O\) 的两个不同点，直线 \(OA\) 和 \(OB\) 的倾斜角分别为 \(\alpha\) 和 \(\beta\)，当 \(\alpha, \beta\) 变化且 \(\alpha + \beta\) 为定值 \(\theta (0 < \theta < \pi)\)，证明直线 \(AB\) 恒过定点，并求出该定点的坐标 \fillinblank{}.
\end{problem}

\begin{problem}
如图，已知长方体 \(ABCD - A_{1}B_{1}C_{1}D_{1}\)，\(AB = 2\)，\(AA_{1} = 1\)，直线 \(BD\) 与平面 \(AA_{1}B_{1}B\) 所成的角为 \(30^{\circ}\)，\(AE\) 垂直 \(BD\) 于 \(E, F\) 为 \(A_{1}B_{1}\) 的中点. (1) 求异面直线 \(AE\) 与 \(BF\) 所成的角; (2) 求平面 \(BDF\) 与平面 \(AA_{1}B\) 所成的二面角 (锐角) 的大小 (3) 求点 A 到平面 BDF 的距离 \fillinblank{}.
\end{problem}

